\chapter*{Conclusions of Part~II}
\addcontentsline{toc}{chapter}{Conclusions of Part~II}
\label{ch:part2-concl}

The four chapters in this Part illustrate the breadth of observational and theoretical tools available for probing the dark sector, and the care required in drawing robust conclusions from current data.

On the dark matter side, Chapter~\ref{ch:minimal-dm-freezein} demonstrated that the freeze-in production mechanism yields a qualitatively different phenomenology when the reheating temperature lies below the dark matter mass. The Boltzmann suppression of production in this regime demands larger portal couplings to match the observed relic abundance, enhancing the dark matter--electron scattering cross section by a factor that scales as $\sqrt{x_\mathrm{rh}}\,e^{x_\mathrm{rh}}$ with $x_\mathrm{rh} = m_\chi/T_\mathrm{rh}$. This result reframes the freeze-in benchmark as an extended region in parameter space rather than a single curve, with a large portion already being probed by current direct detection experiments. In Chapter~\ref{ch:cmb-leptophilic-dm}, we showed that for sub-MeV leptophilic dark matter the leading observational handle is not direct scattering but the loop-induced two-photon decay channel. Using a cosmologically consistent treatment of energy deposition within the \textit{Planck} NPIPE likelihood, we derived CMB bounds that surpass existing terrestrial constraints by up to two orders of magnitude for pseudo-scalar dark matter lighter than $\sim 5\,$keV. These two chapters together highlight the complementarity between cosmological probes and laboratory experiments in mapping the viable parameter space for dark matter.

Moving beyond particle candidates, Chapter~\ref{ch:pbh-memory-burden} subjected the memory-burden hypothesis to a suite of cosmological constraints and found that the proposed light primordial black hole mass window survives only under an extremely restrictive condition. The onset of memory-burden stabilization must occur almost immediately, with the black hole retaining all but a fraction $\lesssim 10^{-10}$ of its initial mass before evaporation is suppressed. Any smooth or gradual transition is ruled out by CMB and BBN data. This result does not invalidate the memory-burden mechanism itself but establishes that any modification to Hawking evaporation must be both extreme and abrupt to evade observational bounds.

Finally, Chapter~\ref{ch:desi-dark-energy} addressed the interpretation of recent DESI DR2 evidence for evolving dark energy. By constructing theory-informed priors from thawing quintessence models motivated by string theory and axion physics and applying them through a normalizing-flow reweighting framework, we showed that the significance of the dark energy signal drops from ${\sim}3\sigma$ under uniform priors to ${\sim}1.3$--$1.8\sigma$ under physically motivated ones, rendering it consistent with a cosmological constant. This finding underscores that even ostensibly uninformative flat priors carry implicit assumptions in phenomenological parameter spaces and that marginal detections are particularly susceptible to prior-induced bias.

Taken together, these results demonstrate the power of combining precision cosmological data with careful theoretical modeling. Whether constraining the couplings of feebly interacting dark matter, stress-testing speculative modifications to black hole physics, or assessing the robustness of potential departures from $\Lambda$CDM, the common lesson is that the interpretation of observational evidence for the dark sector is inseparable from the theoretical framework in which it is analyzed.
